\hnsection{李群概念的扩展}

如今，李理论的活力体现在其应用的多样性（拓扑学、微分几何、算术等）以及平行理论的创立上；在这些理论中，底层流形结构被类似的结构所替代（ \(p\) 进流形、代数簇、概形、形式概形、\(\ldots\) ）。我们这里不讨论这些发展的历史，而只限于第三章涉及的内容：巴拿赫李群和 \(p\) 进李群。

\textbf{(a) 巴拿赫李群}

这里讨论的是“无限维”李群。从局部观点看，欧氏空间中的 0 的邻域被巴拿赫空间中的 0 的邻域所替代。这就是 G. Birkhoff 于 1936 年所做的工作 {[}29 a{]}，由此建立了完备赋范李代数的概念，以及它与定义在巴拿赫空间开集上的“群芽”之间的对应。约在 1950 年，Dynkin 将豪斯多夫公式推广到这一情形，从而完成了这些结果（cf.~supra）。

Birkhoff 和 Dynkin 的定义和结果都是局部的。直到最近，似乎还没有人试图明确建立相应的全局理论，这无疑是由于缺乏应用。†

\textbf{(b) p 进李群}

这类群最早见于 1907 年 Hensel 的工作 {[}19{]}，内容涉及 \(p\) 进解析函数（解析函数由整级数展开定义）。他特别研究了指数和对数；尽管定义它们的级数的先验行为令人惊讶（例如指数级数并不处处收敛），它们的基本函数性质仍然成立，从而在加法群与 \(\mathbf{Q}_p\) 的乘法群之间给出局部同构（或者更一般地，在任意特征为零的完备超度量域的加法群与乘法群之间给出局部同构）。

A. Weil {[}33{]} 和 E. Lutz {[}34{]} 在关于 p 进椭圆曲线（1936）的工作中同样研究了交换群（不过这次是非线性的）。除算术应用外，这里还基于对不变微分形式的积分，构造了群与加法群之间的局部同构。正如 C. Chabauty 不久后指出的，这一方法同样适用于阿贝尔簇；他未作进一步说明就用它证明了“莫德尔猜想”的一个特例 {[}35{]}。

从那时起，很明显，李群的局部理论几乎无需改变就可以应用于 p 进情形。李群--李代数“词典”的基本定理由 Chevalley 的学生 R. Hooke 在 1942 年的学位论文中建立 {[}37{]}；这项工作还包含 E. Cartan 关于实李群闭子群的定理的 p 进类比。

近来，M. Lazard {[}42 b{]} 为定义在 \(\mathbf{Q}_{\mathrm{p}}\) 上的紧解析群发展了更精确的“词典”形式。他证明，紧群 G 上存在 \(p\) 进解析结构，与 G 上某些滤过的存在密切相关，并由此给出各种应用（例如 G 的上同调）。Lazard 的工具之一，是改进 Dynkin 关于 \(p\) 进豪斯多夫级数收敛性的结果 {[}42 a{]}。

